% =====================================================================
%  CHAPTER 8 - CONCLUSION AND FUTURE WORK
% =====================================================================
\clearpage
\chapter{Conclusion and Future Work}

\section{Conclusion}

This project set out to build a durable execution engine that partitions
long-running workflows across several nodes, survives the failure of any of
them without losing progress or repeating side effects, and needs no
infrastructure other than PostgreSQL. The engine records every decision of a
workflow in an append-only history and recovers a run by replaying its code
against that history. Workflows are placed on a fixed set of shards, and
shards on nodes by a consistent-hash ring, so that ownership can move without
moving data. Ownership is held through a lease row per shard, and every
transaction that writes shard state re-checks the shard's fencing token under
a row lock before it commits.

The evaluation on Amazon EKS supports the main claims within the limits set
out in Section~\ref{sec:validity}. Under sixteen valid fault trials of six kinds, including
a 35\,s freeze of an engine process that outlived its lease, no journal
invariant was violated and every accepted run completed; the frozen engine's
attempts to write after it resumed were rejected by fencing in every trial.
After a graceful deletion or a process kill, all of a node's shards were
re-owned within 6\,s. The cost of faults was availability, in the form of
refused starts and longer tail latency while shards were unowned, and never
correctness.

The evaluation also located the limits of the design. On a single 2-vCPU
database host the engine sustained about 131 workflows per second in a
closed loop and about 100 in an open loop, and adding engine replicas did not
change this, because the database was the bottleneck in every configuration.
Much of the database's work was polling: an idle cluster committed about 250
transactions per second, and at the knee about 100 transactions were spent
per completed workflow. Polling also set a latency floor of about 2.2\,s for a
workflow with four activities and a short timer. Under open-loop overload the
system lost shard leases, which turned overload into unavailability. And a
general-purpose sentence encoder proved to be a poor basis for semantic
deduplication on its own, judging most distinct but similar requests to be
duplicates; comparing the deciding fields of a request exactly, and using
similarity only for the remainder, resolved this on the evaluation corpus.

\section{Future Work}

The results point to the following work, in roughly the order of its
expected benefit.

\begin{enumerate}[leftmargin=*, itemsep=2pt]
  \item \emph{Notification instead of polling.} Waking workers with
    PostgreSQL \cd{LISTEN}/\cd{NOTIFY} when tasks and timers become ready, and
    polling only as a fallback, should lower both the idle load and the
    latency floor, and raise capacity on the same database.
  \item \emph{Isolated lease renewal.} Renewing leases on a dedicated
    connection, outside the pool used by workers, should prevent the lease
    loss seen under open-loop overload. Admission control that refuses work
    before the pool is exhausted would complement it.
  \item \emph{Faster adoption of orphaned shards.} Making the adoption delay
    depend on whether the owner answers its health check would shorten
    unavailability after a freeze or a drain without causing churn after
    short pauses.
  \item \emph{Safe migrations.} Migrations run under an advisory lock, and
    skipped when the schema is current, would remove the start-up deadlock.
  \item \emph{Scaling and availability of the database.} Repeating the
    evaluation with a larger database host, and with a synchronous standby
    or a managed highly available database, is needed to measure horizontal
    scaling of the engine itself and to test failover of the database.
  \item \emph{History retention.} Histories of completed runs should be
    archived or compacted.
  \item \emph{Semantic deduplication on real traffic.} The hybrid comparator
    should be evaluated on agent traces labelled by people, with thresholds
    chosen on a held-out split, and the escalation path should be studied as
    a human-in-the-loop interaction.
  \item \emph{Wider fault coverage.} Network partitions, packet loss,
    instance termination and rolling upgrades under load remain untested.
\end{enumerate}
